% Part: first-order-logic
% Chapter: models-theories
% Section: introduction

\documentclass[../../../include/open-logic-section]{subfiles}

\begin{document}

\olfileid{fol}{mat}{int}
\olsection{परिचय}

\begin{explain}
स्वयंसिद्धकीय पद्धतीचा विकास हा विज्ञानाच्या इतिहासातील एक महत्त्वपूर्ण
पराक्रम आहे आणि गणिताच्या इतिहासात त्याला विशेष महत्त्व आहे. एखाद्या
क्षेत्राचा स्वयंसिद्धकीय विकास करताना अनेक प्रश्न स्पष्ट करावे लागतात: हे
क्षेत्र कशाविषयी आहे? सर्वांत मूलभूत संकल्पना कोणत्या? त्यांचे परस्पर संबंध
कसे आहेत? क्षेत्रातील सर्व संकल्पना या मूलभूत संकल्पनांच्या साहाय्याने
परिभाषित करता येतात का? या संकल्पना कोणते नियम पाळतात आणि त्यांनी कोणते
नियम पाळलेच पाहिजेत?

स्वयंसिद्धकीय पद्धत आणि तर्कशास्त्र जणू एकमेकांसाठीच निर्माण झाले आहेत.
आकारिक तर्कशास्त्र स्वयंसिद्धकीय उपपत्ती मांडण्यासाठी, उपपत्तीच्या
स्वयंसिद्धकांपासून प्रमेये नेमक्या निर्दिष्ट रीतीने सिद्ध करण्यासाठी आणि
स्वयंसिद्धकांची पूर्ति करणाऱ्या सर्व प्रणालींच्या गुणधर्मांचा सुव्यवस्थित
अभ्यास करण्यासाठी साधने पुरवते.
\end{explain}

\begin{defn}
!!{sentence}s चा संच~$\Gamma$ \emph{बंद} असतो तेव्हा आणि केवळ तेव्हाच,
जेव्हा $\Gamma \Entails !A$ असल्यास $!A \in \Gamma$ असते.
!!{sentence}s च्या संच~$\Gamma$ चे \emph{संवरण} म्हणजे
$\Setabs{!A}{\Gamma \Entails !A}$ होय.

जर~$\Gamma$ हे~$\Delta$ चे संवरण असेल, तर~$\Gamma$ हा वाक्यांच्या
संच~$\Delta$ ने \emph{स्वयंसिद्धकांद्वारे निरूपित} केला आहे असे म्हणतो.
\end{defn}

\begin{explain}
एखादी स्वयंसिद्धकीय उपपत्ती म्हणजे तिच्या स्वयंसिद्धकांच्या संच~$\Delta$
ने निरूपित केलेल्या वाक्यांचा संच, असे आपण मानू शकतो. दुसऱ्या शब्दांत,
आपल्याला अभ्यासायच्या स्वयंसिद्धकीय रीतीने विकसित केलेल्या शास्त्रातील
मूलभूत संकल्पना दर्शवणारी अतार्किक चिन्हे असलेली प्रथम-क्रम भाषा आणि त्या
शास्त्राचे मूलभूत नियम व्यक्त करणारा !!{sentence}s चा संच दिलेला असेल,
तर स्वयंसिद्धकांपासून तार्किकरीत्या निष्पन्न होणारी या भाषेतील सर्व
!!{sentence}s उपपत्तीचे प्रतिनिधित्व करतात असे आपण मानू शकतो. याची व्याप्ती,
एकच आदिम संकल्पना आणि सोपी स्वयंसिद्धके असलेल्या अंशतः क्रमांच्या
उपपत्तीसारख्या साध्या उदाहरणांपासून न्यूटनीय यामिकीसारख्या गुंतागुंतीच्या
उपपत्तींपर्यंत आहे.

स्वयंसिद्धकीय पद्धतीकडे पाहण्याचा हा आकारिक दृष्टिकोन इतका महत्त्वाचा
ठरण्यामागील मुख्य तार्किक तथ्ये पुढीलप्रमाणे आहेत. $\Gamma$ ही एखाद्या
उपपत्तीची स्वयंसिद्धक-प्रणाली, म्हणजे वाक्यांचा संच, आहे असे समजा.
\begin{enumerate}
\item एखादी स्वयंसिद्धक-प्रणाली !!{structure}s चा अभिप्रेत वर्ग नेमका कधी
  पकडते, हे आपण काटेकोरपणे सांगू शकतो. म्हणजे, आपल्याला !!{structure}s च्या
  एखाद्या विशिष्ट वर्गात रस असेल, तर स्वयंसिद्धक-प्रणाली~$\Gamma$ तो वर्ग
  तेव्हा आणि केवळ तेव्हाच यशस्वीपणे पकडते, जेव्हा त्या !!{structure}s म्हणजे
  नेमकी ती~$\Struct M$ असतात ज्यांच्यासाठी $\Sat{M}{\Gamma}$.
\item या बाबतीत आपल्याला अपयश येऊ शकते, कारण काही~$\Struct M$ साठी
  $\Sat{M}{\Gamma}$ असते, पण~$\Struct M$ ही आपण अभिप्रेत धरलेल्या
  !!{structure}s पैकी नसते. त्यामुळे~$\Struct M$ मध्ये सत्य नसलेली
  स्वयंसिद्धके आपल्याला जोडावी लागू शकतात.
\item निदान~$\Gamma$ ही सर्व अभिप्रेत !!{structure}s मध्ये सत्य आहे एवढ्या
  बाबतीत आपण यशस्वी झालो, तर $\Gamma \Entails !A$ असताना वाक्य~$!A$ हे
  सर्व अभिप्रेत !!{structure}s मध्ये सत्य असते. म्हणून, वाक्ये सर्व
  अभिप्रेत !!{structure}s मध्ये सत्य आहेत हे दाखवण्यासाठी, ती
  स्वयंसिद्धकांपासून तार्किकरीत्या निष्पन्न होतात एवढे दाखवून आपण तार्किक
  साधने (उदाहरणार्थ, !!{derivation} पद्धती) वापरू शकतो.
\item कधीकधी आपल्या मनात अभिप्रेत !!{structure}s नसतात; त्याऐवजी आपण
  स्वयंसिद्धकांपासूनच सुरुवात करतो: काही विशिष्ट नियमांची पूर्ति करावी अशा
  काही आदिम संकल्पना आपण घेतो आणि ते नियम स्वयंसिद्धक-प्रणालीत संकेतबद्ध
  करतो. स्वयंसिद्धके परस्परविरोधी नाहीत, ही गोष्ट आपल्याला लगेच पडताळून
  पाहायची असते. ती परस्परविरोधी असतील, तर त्या नियमांचे पालन करणाऱ्या
  संकल्पना असूच शकत नाहीत आणि आपण विसंगत उपपत्ती मांडण्याचा प्रयत्न केलेला
  असतो. $\Gamma$ चे प्रतिमान शोधून असे घडत नाही हे आपण पडताळू शकतो. आणि
  आपल्या उपपत्तीची प्रतिमाने असतील, तर त्यांचा तपास करण्यासाठी तसेच प्रतिमाने
  रचण्यासाठी आपण तार्किक पद्धती वापरू शकतो.
\item स्वयंसिद्धकांचे स्वातंत्र्य हाही एक महत्त्वाचा प्रश्न आहे. एखादे
  स्वयंसिद्धक प्रत्यक्षात इतरांचा परिणाम असल्यामुळे अनावश्यक असू शकते.
  $\Gamma$ मधील स्वयंसिद्धक~$!A$ अनावश्यक आहे हे आपण
  $\Gamma \setminus \{!A\} \Entails !A$ सिद्ध करून दाखवू शकतो. तसेच,
  $(\Gamma \setminus \{!A\}) \cup \{\lnot !A\}$ पूर्ततायोग्य आहे हे
  दाखवून ते स्वयंसिद्धक अनावश्यक नाही हेही सिद्ध करू शकतो. उदाहरणार्थ,
  भूमितीचे समांतर गृहीतक इतर स्वयंसिद्धकांपासून स्वतंत्र आहे हे अशाच प्रकारे
  दाखवले गेले.
\item उपपत्तीतील संकल्पनांची परिभाष्यता हाही आणखी एक महत्त्वाचा प्रश्न आहे:
  भाषेच्या निवडीवर उपपत्तीची प्रतिमाने कशाची बनलेली असतात हे ठरते. पण
  उपपत्तीचा प्रत्येक पैलू तिच्या प्रतिमानांत स्वतंत्रपणे दर्शवावाच लागतो असे
  नाही. उदाहरणार्थ, प्रत्येक क्रमसंबंध~$\le$ त्याच्याशी संबंधित काटेकोर
  क्रमसंबंध~$<$ ठरवतो---एक दिला असेल, तर दुसरा परिभाषित करता येतो. म्हणून
  असा क्रम असलेल्या उपपत्तीच्या प्रतिमानात त्याच्याशी संबंधित काटेकोर
  क्रमसंबंध \emph{वेगळ्याने} असणे आवश्यक नाही. सर्वसाधारणपणे, एक संबंध
  इतरांच्या आधारे नेमका कधी परिभाषित करता येतो? एखादा संबंध इतरांच्या
  आधारे परिभाषित करणे कधी अशक्य असते (आणि म्हणून तो भाषेच्या आदिम
  संकल्पनांत जोडावा लागतो)?
\end{enumerate}
\end{explain}

\end{document}
